%==============================================================
% DISCUSSION + CONCLUSION -- journal version.
%
% Defines sec:limitations, referenced from the Related Work section.
%
% The shift interpretation here is the OPPOSITE of the conference version's.
% That version argued the shift penalty was mostly a calibration effect; the
% artefact control shows that was an artefact of the uncontrolled encoding
% (Section IV-D), and the behaviour of the two interventions points the other
% way. Section A below gives the reasoning and the numbers behind it.
%==============================================================

\section{Discussion}

\subsection{One Benchmark, Two Tasks}
\label{sec:twotasks}

The per-split decomposition is the result with consequences beyond this
system. BanFakeMM distributes five adversarial sets under one name and one
labelling convention, and it is natural to read them as five difficulty levels
of a single task. They are not. Two of them edit or synthesise the photograph;
the other three rewrite the article and ship the original photograph unchanged,
byte for byte. This is, we presume, by design: attacks that alter the text and
keep the photograph are the natural way to build cross-modal inconsistency, and
they are well posed for a detector that reads both. The point is therefore not
that the benchmark is flawed but that an image-only evaluation falls outside its
intended use on those splits. A model reading only pixels is answering a
different question in each group, and in the second group it is answering a
question whose answer is absent from its input.

This matters because the failure is silent. The model does not crash or abstain
on LinguisticAttack; it returns $49.5$ percent accuracy, which looks like a
weak result rather than an undefined one. An author who reports an image-model
average across all five splits produces a number that is two parts measurement
and three parts noise, and nothing in the accuracy figure reveals the
difference. The check that does reveal it is cheap---hash the manipulated
images against the authentic corpus---and we would argue it belongs in the
reporting of any multi-attack benchmark, alongside the split sizes.

For this system the consequence is a scope statement. \sysname{} is
image-derived throughout, so we evaluate it where the image carries the
manipulation and say plainly that it is blind elsewhere. A detector that
covered the other three attacks would need the article text as an input, which
the benchmark distributes and which our pipeline does not read. That is the
clearest direction for extending this work, and it is a different system rather
than a tuning of this one.

\subsection{Where the Shift Penalty Actually Lives}
\label{sec:shiftdiscussion}

Under the artefact control, the zero-shot model falls from $94.35$ percent
same-distribution to $84.9$ on the sealed shift partition. Two interventions
are available, and they behave very differently.

Moving the operating point---selecting a threshold on a disjoint validation
partition instead of using the default---buys $+0.5$ points
(95 percent CI $[-3.6, +4.6]$, McNemar $p = 1.000$). Changing the model, by
adding 300 target-style training images, buys $+3.8$ accuracy points and
$+0.027$ AUROC, with intervals excluding zero across eight seeds. A threshold
move cannot change a ranking; adaptation can, and did. The shift penalty is
therefore predominantly representational: the zero-shot model does not merely
place its cut-off badly, it orders the target-style images less well, and only
retraining repairs that.

\emph{This reverses the interpretation we published in the conference version
of this work}, which attributed almost the entire shift penalty to calibration.
That conclusion rested on an uncontrolled measurement in which thresholding
appeared to recover $+6.1$ points. Once the storage artefact is removed, the
same comparison yields $+0.5$. The artefact had inflated the apparent value of
recalibration because it gave the zero-shot model a spurious, well-separated
signal whose decision boundary was simply mis-placed---exactly the situation in
which moving a threshold looks powerful.

We note the boundary of the shift being studied. AdversarialEdit and
AdversarialGen differ in how the manipulated class is produced, not in image
source or capture conditions, and both draw their authentic images from the
same corpus. A shift that also changed provenance or compression could damage
the representation further, and nothing here bounds that case.

\subsection{What a Fixed Threshold Hides}
\label{sec:metricdiscussion}

The adaptation result is a methodological warning as much as an empirical one.
The same intervention, measured on the same images, reads as $-1.5$ points on
194 images, $+0.7$ or $+3.8$ points on 1{,}308 images from a single fit
(original or re-encoded), or $+1.4$ and $+4.0$ points averaged over eight
seeds, depending on the partition size, the encoding, the seed and the operating
point---and as a stable $+0.023$ to $+0.027$ AUROC regardless of any of them. Accuracy at a fixed threshold is a projection
of a model's behaviour onto one point of its ROC curve, and an intervention
that improves the curve elsewhere is invisible to it.

This is not an argument against reporting accuracy, which is what a deployed
system delivers. It is an argument against using accuracy at a single
threshold, on a small partition, from a single fit, to support a \emph{null}.
A null claim asserts that an effect is absent; to carry that weight the
measurement must be sensitive enough to have found the effect had it been
there. The conference version's interval, $[-4.6, +1.5]$ (converged: $[-5.2,
+1.5]$), was wide enough to
contain the $+1.4$ we later measured on $6.7\times$ more data. The honest
reading at the time would have been that the study could not resolve the
question, which is the reading we now recommend for any accuracy-based null on
a partition of that size.

\subsection{Limitations and Future Directions}
\label{sec:limitations}

Ground-truth labels derive from BanFakeMM's deterministic construction, so the
demonstrated competence is specific to detecting algorithmic manipulation, not
to adjudicating whether a news item's claim is true. Establishing that a
photograph was edited says nothing about the story attached to it.

Our normalisation removes the storage artefact but not a generation-pipeline
signature that survives resampling. If every manipulated image carries the same
synthesis fingerprint, a visual model may be recognising the generator rather
than the manipulation, and a benchmark that defines its labels by the
construction pipeline cannot distinguish the two. Settling this needs
manipulated images from several independent generators, which BanFakeMM does
not provide.

The shift evaluation, though sealed, screened at both the image and article
level, and artefact-controlled, draws on a second split of the same benchmark.
The metadata check of Section~\ref{sec:artefact} was repeated on one other
corpus, but the detector itself was not.
Generalization to manipulation styles and image sources outside BanFakeMM,
including organically collected social media imagery, remains open and is the
limitation we would most like to remove.

Two limits remain on the caption-pathway comparison. Under the artefact control
it was repeated only for the same-distribution setting
(Section~\ref{sec:fusion}); the shift-partition comparison, where the fused
model's calibration deficit appeared, was not, because it would need captions for
the shift images as well. And the CISF hyperparameters $k$ and $\alpha$ were fixed
by design rather than search; a post hoc grid over $k \in \{3, 5, 10\}$ and
$\alpha \in \{0.5, 0.65, 0.8\}$ moves fused accuracy only between $94.30$ and
$94.68$ percent, well inside the $1.0$ to $1.3$ point fold-to-fold spread, so the
choice does not drive any conclusion here.

Our comparison against MultiBanFakeDetect \cite{multibanfake} is of feature
families under one protocol, not of trained systems: its text branch consumes
the article's own content, which this pipeline does not read. A comparison that
gave both systems the inputs they were designed for would be more informative
than either, and would also reach the three attack types we are blind to.

Finally, pairing this detector with retrieval over a Bangla fact-check corpus
is the obvious extension, and the one that would address the three attack types
we are blind to. Evaluating such a stage needs a labelled claim set with
adjudicated verdicts, retrieval recall@$k$ against gold evidence, and an
abstention analysis. A preliminary measurement suggests it is harder than it
looks: the multilingual encoder we tested does not separate Bangla fact-check
headlines well enough for reliable retrieval, as they share a fixed rhetorical
template. Architecture-level alternatives to our data-centric mitigations, such
as CLIPood \cite{clipood} and domain-specific prompt tuning \cite{dpod}, were
not evaluated, and extending the pipeline to video would need a different
approach.

\section{Conclusion}

We set out to build a detector for manipulated Bangla news imagery and found
that the harder problem was establishing what any number measured on a
constructed benchmark actually means.

Three properties of BanFakeMM had to be dealt with before its results could be
interpreted. Its manipulated images are separable from authentic ones by file
metadata alone, at $100$ percent accuracy and without reading a pixel; we
re-encode every image through one pipeline, which removes the artefact at a
cost of $1.03$ accuracy points and establishes that the remaining signal is in
the content. A weaker form of the same shortcut, $77$ to $83$ percent from
metadata alone, appears in a second Bangla corpus, so it is not specific to this
benchmark. Its five attack types do not pose one task: in three of them every
manipulated image is byte-identical to an authentic one, so an image-only model
is undecidable there, and we scope our claims to the two splits where the
manipulation reaches the pixels. And its images recur---both as near-duplicates
and as repeated source articles---so every partition is screened on content and
the residual article-level overlap is measured rather than assumed away.

Within that scope, the findings are these. Visual features carry the
discriminative signal: CLIP alone reaches $95.38$ percent under duplicate-aware
grouped cross-validation, $94.35$ with the artefact removed, while every
text-only configuration saturates below $78$. OCR is near enough to chance to
remove, and the caption-fusion pathway returns nothing measurable for its
computational cost. Transfer to an unseen manipulation style costs roughly nine
points of accuracy. Supervised adaptation on 300 target-style images recovers a
substantial part of that, improving ranking by $+0.025$ AUROC with intervals
excluding zero across eight seeds and both encodings.

Two of these conclusions correct our own earlier published claims about this
system: that the shift penalty was mostly a calibration effect, and that
supervised adaptation added nothing. Both errors had the same cause---a
measurement that was too small, too threshold-dependent, and taken on
uncontrolled data---and both were invisible from inside the result that
produced them. We report the corrections with the measurements that force them,
because a field that cannot revise its own numbers in public accumulates them
instead.

The gap between a high same-distribution figure and reliable behaviour outside
the training distribution is real. What this work adds is that a sizeable part
of such a gap, and sometimes its entire apparent explanation, can be a property
of the benchmark rather than of the model---and that distinguishing the two
costs a hash, a re-encode, and the discipline to measure an effect with a
statistic sensitive enough to find it.
